\begin{ouverture}[probabilites]{21}{Premiers pas vers le hasard}{Partie 5 · Modéliser et organiser}
{Expériences aléatoires · Issues · Échelle des chances}
\objectif{Reconnaître une expérience aléatoire et lister ses issues.}
\objectif{Employer le vocabulaire : certain, possible, impossible, probable.}
\objectif{Comparer les chances de réalisation d’évènements.}
\prerequis{Fractions simples et dénombrement.}
\end{ouverture}
\pagestyle{otMain}\essentiel\fontsize{11}{15}\selectfont
\phantomsection\label{cours-21-1}
\essSection{1}{Le vocabulaire du hasard}
\begin{definition}Une expérience est aléatoire quand son résultat ne peut pas être prévu avant de la réaliser : lancer un dé, tirer une boule d’un sac sans regarder, piocher une carte. Les résultats possibles s’appellent les issues.\end{definition}
\begin{exemple}Avec un dé ordinaire : « obtenir 7 » est impossible ; « obtenir un nombre de 1 à 6 » est certain ; « obtenir 2 » est possible, sans certitude. Entre impossible et certain, tous les degrés existent.\end{exemple}
\label{fig-cours-21-1-cours-echelle}\figureGeo{cours-echelle}

\begin{savoirfaire}[Une chance sur deux]« Une chance sur deux » décrit un évènement qui a autant de chances de se réaliser que de ne pas se réaliser : obtenir pile avec une pièce, ou tirer une carte rouge dans un jeu qui compte autant de rouges que de noires.\end{savoirfaire}
\clearpage
\phantomsection\label{cours-21-2}
\essSection{2}{Compter les issues}
\begin{savoirfaire}[Recenser les issues]Lister toutes les issues avant de raisonner : une pièce en a 2, un dé en a 6, une urne de 4 boules en a 4. Un évènement est d’autant plus probable qu’il est réalisé par davantage d’issues, quand toutes les issues ont les mêmes chances.\end{savoirfaire}
\begin{definition}Quand toutes les issues ont les mêmes chances de sortir, la situation est équiprobable : dé équilibré, boules indiscernables au toucher, cartes bien mélangées.\end{definition}
\begin{attention}Deux issues possibles ne donnent pas toujours une chance sur deux ! Avec une urne de 1 boule rouge et 3 bleues, « rouge » et « bleu » sont deux résultats possibles, mais le bleu a trois fois plus d’issues pour lui.\end{attention}
\begin{exemple}Obtenir 6 trois fois de suite n’assure rien pour le lancer suivant : le résultat reste imprévisible, c’est justement une expérience aléatoire.\end{exemple}

\clearpage\rubrique{Découverte}
\begin{blocLivre}{ex-21-1}\begin{exercice}[1]{Aléatoire ou pas ?}
Dire si l’expérience est aléatoire : lancer un dé et lire la face ; faire bouillir de l’eau et observer si elle chauffe ; tirer une carte d’un jeu mélangé.
\end{exercice}\end{blocLivre}
\begin{blocLivre}{ex-21-2}\begin{exercice}[1]{Le bon mot}
Avec un dé ordinaire, qualifier chaque évènement : « obtenir 7 » ; « obtenir un nombre de 1 à 6 » ; « obtenir 2 ».
\end{exercice}\end{blocLivre}
\begin{blocLivre}{ex-21-3}\begin{exercice}[1]{Lister les issues}
Lister toutes les issues : d’un lancer de pièce ; d’un lancer de dé ; du tirage d’une boule dans un sac contenant une boule rouge, une verte et une bleue.
\end{exercice}\end{blocLivre}
\begin{blocLivre}{ex-21-4}\begin{exercice}[1]{Une chance sur deux}
Dire si l’évènement a une chance sur deux de se réaliser : obtenir pile avec une pièce ; obtenir un nombre pair avec un dé ; tirer une carte rouge d’un jeu de 32 cartes bien mélangé.
\end{exercice}\end{blocLivre}
\begin{blocLivre}{ex-21-5}\begin{exercice}[1]{Sur l’échelle des chances}
Placer chaque évènement sur l’échelle, d’impossible à certain : « janvier suivra décembre » ; « obtenir 3 avec un dé » ; « tirer une boule jaune d’un sac qui n’en contient pas ».
\label{fig-ex-21-5}\figureGeo{pro-5}
\end{exercice}\end{blocLivre}

\rubrique{Entraînement}
\begin{blocLivre}{ex-21-6}\begin{exercice}[2]{Plus ou moins probable}
Avec un dé équilibré, quel évènement est le plus probable : « obtenir 2 » ou « obtenir un nombre pair » ? Justifier en comptant les issues.
\label{fig-ex-21-6}\figureGeo{pro-6}
\end{exercice}\end{blocLivre}
\begin{blocLivre}{ex-21-7}\begin{exercice}[2]{L’urne trompeuse}
Un sac contient 1 boule rouge et 3 boules bleues, indiscernables au toucher. Maé dit : « Rouge ou bleu, il y a deux résultats : j’ai une chance sur deux de tirer rouge. » Corriger.
\label{fig-ex-21-7}\figureGeo{pro-7}
\end{exercice}\end{blocLivre}
\begin{blocLivre}{ex-21-8}\begin{exercice}[2]{Équiprobable ou non ?}
Dire si toutes les issues ont les mêmes chances : un dé équilibré ; la roue dessinée, aux secteurs de tailles différentes ; une pièce truquée qui retombe presque toujours sur face.
\label{fig-ex-21-8}\figureGeo{pro-8}
\end{exercice}\end{blocLivre}
\begin{blocLivre}{ex-21-9}\begin{exercice}[2]{Le sac de la kermesse}
Un sac contient 5 jetons rouges et 2 jetons bleus. L’évènement « tirer un rouge » est-il probable ou peu probable ? Et « tirer un bleu » ?
\label{fig-ex-21-9}\figureGeo{pro-9}
\end{exercice}\end{blocLivre}
\begin{blocLivre}{ex-21-10}\begin{exercice}[2]{Les jetons numérotés}
Un sac contient huit jetons numérotés de 1 à 8. Quel évènement est le plus probable : « tirer un numéro plus grand que 6 » ou « tirer un numéro plus petit que 4 » ?
\end{exercice}\end{blocLivre}

\rubrique{Problèmes}
\begin{blocLivre}{ex-21-11}\begin{exercice}[3]{Le jeu de 32 cartes}
Dans un jeu de 32 cartes bien mélangé, quel évènement est le plus probable : « tirer un cœur » ou « tirer une carte rouge » ? Justifier.
\end{exercice}\end{blocLivre}
\begin{blocLivre}{ex-21-12}\begin{exercice}[3]{Gagner ou perdre}
À une tombola, Jules dit : « Je gagne ou je perds, donc j’ai une chance sur deux de gagner. » Que lui répondre, sachant que 500 billets sont vendus et que 10 gagnent ?
\end{exercice}\end{blocLivre}
\begin{blocLivre}{ex-21-13}\begin{exercice}[3]{La roue de la kermesse}
La roue de la kermesse porte 8 secteurs identiques : 4 marqués « perdu », 3 « bonbon » et 1 « peluche ». Quel lot est le plus probable ? La peluche a-t-elle une chance sur deux de sortir ?
\label{fig-ex-21-13}\figureGeo{pro-13}
\end{exercice}\end{blocLivre}
\begin{blocLivre}{ex-21-14}\begin{exercice}[3]{La série de six}
Nora a obtenu 6 trois fois de suite avec un dé équilibré. Elle affirme : « Je suis sûre d’avoir encore 6 au prochain lancer. » Puis elle se ravise : « Non, le 6 ne peut plus sortir ! » Que penser de ces deux phrases ?
\end{exercice}\end{blocLivre}
\begin{blocLivre}{ex-21-15}\begin{exercice}[3]{Classer quatre évènements}
On tire une boule d’un sac contenant 2 rouges, 3 vertes et 5 noires. Classer du moins probable au plus probable : « tirer une rouge », « tirer une verte », « tirer une noire », « tirer une boule rouge, verte ou noire ».
\label{fig-ex-21-15}\figureGeo{pro-15}
\end{exercice}\end{blocLivre}

\clearpage
